% Part:sets-functions-relations
% Chapter: size-of-sets
% Section: pairing-alt

\documentclass[../../../include/open-logic-section]{subfiles}

\begin{document}

\olfileid[fa-IR]{sfr}{siz}{pai-alt}

\olsection{یک تابع جفت‌سازی جایگزین}

\begin{explain}
برشماری‌های دیگری از $\Nat^2$ وجود دارند که یافتن وارونشان را آسان‌تر
می‌کنند. یکی از آن‌ها را در اینجا می‌بینیم. به‌جای آنکه برشماری را در
یک آرایه مجسم کنید، با فهرستی از اعداد صحیح مثبت آغاز کنید که با
جایگاه‌هایی (در ابتدا) تهی متناظرند. تصور کنید این جایگاه‌ها را به‌ترتیب
و به‌شکل زیر با زوج‌های $\tuple{n,m}$ پر می‌کنیم. از زوج‌هایی آغاز کنید
که در مؤلفهٔ نخستشان~$0$ دارند (یعنی زوج‌های $\tuple{0,m}$): نخستین زوج
(یعنی $\tuple{0,0}$) را در نخستین جایگاه تهی بگذارید، سپس یک جایگاه تهی
را رد کنید، دومین زوج (یعنی $\tuple{0,2}$) را در جایگاه تهی بعدی بگذارید،
دوباره یکی را رد کنید، و به همین ترتیب ادامه دهید. آغازِ (ناکاملِ)
برشماری ما اکنون به‌شکل زیر است:
\[\small
\begin{array}{@{}c c c c c c c c c c c@{}}
\mathbf 1 & \mathbf 2 & \mathbf 3 & \mathbf 4 & \mathbf 5 & \mathbf 6 & \mathbf 7 & \mathbf 8 & \mathbf 9 & \mathbf{10} & \dots \\ \\
\tuple{0,0} &  & \tuple{0,1} &  & \tuple{0,2} &  & \tuple{0,3} & & \tuple{0,4} &  & \dots \\
\end{array}
\]
این کار را برای جایگاه‌هایی که هنوز خالی مانده‌اند، با زوج‌های
$\tuple{1,m}$ تکرار کنید و باز هم یک جایگاه خالی در میان را رد کنید:
\[\small
\begin{array}{@{}c c c c c c c c c c c@{}}
\mathbf 1 & \mathbf 2 & \mathbf 3 & \mathbf 4 & \mathbf 5 & \mathbf 6 & \mathbf 7 & \mathbf 8 & \mathbf 9 & \mathbf{10} & \dots \\ \\
\tuple{0,0} & \tuple{1,0} & \tuple{0,1} &  & \tuple{0,2} & \tuple{1,1} &
\tuple{0,3} & & \tuple{0,4} &  \tuple{1,2} & \dots \\
\end{array}
\]
زوج‌های $\tuple{2,m}$، $\tuple{2,m}$ و جز آن را نیز به همین شیوه وارد
کنید. بنابراین برشماری کامل‌شدهٔ ما چنین آغاز می‌شود:
\[\small
\begin{array}{@{}cc c c c c c c c c c@{}}
\mathbf 1 & \mathbf 2 & \mathbf 3 & \mathbf 4 & \mathbf 5 & \mathbf 6 & \mathbf 7 & \mathbf 8 & \mathbf 9 & \mathbf{10} & \dots \\ \\
\tuple{0,0} & \tuple{1,0} & \tuple{0,1} & \tuple{2,0}  & \tuple{0,2} &
\tuple{1,1} & \tuple{0,3} & \tuple{3,0}  & \tuple{0,4} &  \tuple{1,2} & \dots \\
\end{array}
\]
اگر خانه‌های آرایهٔ بالا را مطابق این برشماری شماره‌گذاری کنیم، خط
زیگزاگ مرتبی نخواهیم یافت، بلکه آرایش زیر را خواهیم داشت:
\[
\begin{array}{ c | c | c | c | c | c | c | c }
& \mathbf 0 & \mathbf 1 & \mathbf 2 & \mathbf 3 & \mathbf 4 & \mathbf 5 & \dots \\
\hline
\mathbf 0 & 1 & 3 & 5 & 7 & 9 & 11 & \dots \\
\hline
\mathbf 1 & 2 & 6 & 10 & 14 & 18 & \dots & \dots \\
\hline
\mathbf 2 & 4 & 12 & 20 & 28 & \dots & \dots & \dots \\
\hline
\mathbf 3 & 8 & 24 & 40 & \dots & \dots & \dots & \dots \\
\hline
\mathbf 4 & 16 & 48 & \dots & \dots & \dots & \dots & \dots \\
\hline
\mathbf 5 & 32 & \dots & \dots & \dots & \dots & \dots & \dots \\
\hline
\vdots & \vdots & \vdots & \vdots & \vdots & \vdots & \vdots & \ddots\\
\end{array}
\]
می‌بینیم که زوج‌های سطر~$0$ در جایگاه‌های فردشمارهٔ برشماری ما
قرار دارند؛ یعنی زوج $\tuple{0,m}$ در جایگاه $2m+1$ است. زوج‌های سطر
دوم، یعنی $\tuple{1,m}$، در جایگاه‌هایی قرار دارند که شمارهٔ آن‌ها دو
برابر یک عدد فرد، و به‌طور مشخص $2 \cdot (2m+1)$ است. زوج‌های سطر سوم،
یعنی $\tuple{2,m}$، در جایگاه‌هایی قرار دارند که شمارهٔ آن‌ها چهار برابر
یک عدد فرد، یعنی $4 \cdot (2m+1)$ است؛ و به همین ترتیب. عامل‌های
$(2m+1)$ برای هر سطر، یعنی $1$، $2$، $4$، $8$، \dots، دقیقاً توان‌های~$2$
هستند: $1= 2^0$، $2 = 2^1$، $4 = 2^2$، $8 = 2^3$، \dots\@ در واقع،
نمای مربوط همواره مؤلفهٔ نخست زوج مورد بحث است. پس برای زوج
$\tuple{n,m}$، عامل برابر $2^n$ است. این امر فرمول کلی
$2^n \cdot (2m+1)$ را به ما می‌دهد. بااین‌حال، این نگاشتی از زوج‌ها به
اعداد صحیح \emph{مثبت} است؛ یعنی $\tuple{0,0}$ در جایگاه~$1$ قرار دارد.
اگر بخواهیم از جایگاه~$0$ آغاز کنیم، باید~$1$ را از حاصل کم کنیم. پس داریم:
\end{explain}

\begin{ex}
تابع $h\colon \Nat^2 \to \Nat$ که به‌صورت زیر داده می‌شود
\[
h(n,m) = 2^n (2m+1) - 1
\]
یک تابع جفت‌سازی برای مجموعهٔ زوج‌های اعداد طبیعی~$\Nat^2$ است.
\end{ex}

\begin{explain}
بر این اساس، در دومین برشماریِ $\Nat^2$، زوج $\tuple{0,0}$ دارای کد
$h(0,0) = 2^0(2\cdot 0+1) - 1 = 0$ است؛ زوج $\tuple{1,2}$ دارای کد
$2^{1} \cdot (2 \cdot 2 + 1) - 1 = 2
\cdot 5 - 1 = 9$ است؛ زوج $\tuple{2,6}$ دارای کد $2^{2} \cdot (2
\cdot 6 + 1) - 1 = 51$ است.
\end{explain}

گاهی کافی است زوج‌های اعداد طبیعی~$\Nat^2$ را کدگذاری کنیم، بی‌آنکه
لازم باشد کدگذاری پوشا باشد. چنین کدگذاری‌هایی وارون‌هایی دارند که فقط
توابع جزئی‌اند.

\begin{ex}
تابع $j\colon \Nat^2 \to \Nat^+$ که به‌صورت زیر داده می‌شود
\[
j(n,m) = 2^n3^m
\]
تابعی !!a{injective} به‌صورت $\Nat^2 \to \Nat$ است.
\end{ex}

\end{document}
